% ═══════════════════════════════════════════════════════════════════════
\section{Synthèse}
% ═══════════════════════════════════════════════════════════════════════

\begin{frame}{Ce qu'il faut retenir}
  \small
  \begin{itemize}
    \item Probabilités totales : « multiplier le long d'une branche,
      additionner les branches » --- une ligne d'un produit matriciel.
    \item Chaîne de Markov : l'avenir ne dépend que du présent ; la loi
      évolue par $\vec x_{n+1}=M\vec x_n$, avec $M$ stochastique (colonnes
      de somme $1$).
    \item L'équilibre $\vec q=M\vec q$ : un vecteur propre de valeur
      propre $1$, et un point fixe ; limite des probabilités et proportion
      du temps passé dans chaque état.
    \item Si $M$ est régulière : un unique équilibre, atteint depuis
      n'importe quel départ ; la vitesse dépend de $|\lambda_2|$.
    \item PageRank : corriger les pages sans lien ($B$), puis téléporter
      ($G=\beta B+(1-\beta)U$).
    \item Méthode des puissances : $G$ est une contraction de rapport
      $\beta$ pour $d_1$ ; Banach garantit la convergence et borne
      l'erreur.
  \end{itemize}
\end{frame}

\begin{frame}{Un même point fixe, trois chapitres}
  \small
  \begin{center}
  \begin{tabular}{lll}
    \toprule
    Chapitre & Espace complet & Point fixe \\
    \midrule
    2 & $I\subset\R$ fermé & le zéro cherché par Newton \\
    3 & $\big(\Comp{\R^2},\dH\big)$ & une fractale \\
    4 & $(\Delta,d_1)$ & le PageRank du Web \\
    \bottomrule
  \end{tabular}
  \end{center}
  \vspace{6pt}
  À l'origine de l'une des plus grandes entreprises du monde : un projet
  de doctorat dont le cœur est le calcul d'un vecteur propre associé à la
  valeur propre $1$.
\end{frame}

\begin{frame}{Et maintenant}
  \begin{itemize}
    \item Les démonstrations laissées « au tableau » aujourd'hui sont
      toutes dans les notes de cours, avec tous les détails.
    \item Les exercices du chapitre (avec solutions complètes) couvrent
      chacun des résultats vus : à travailler avant le prochain cours.
    \item Les programmes Python des notes (météo, simulation, mini-webs,
      méthode des puissances) sont à exécuter et à modifier.
    \item Les rappels d'algèbre linéaire (valeurs propres,
      diagonalisation) sont à l'annexe D.
  \end{itemize}
\end{frame}
